
As in the previous section, we begin by introducing refinements of the quadruple sums, along with arbitrary linear terms in the exponent of $q$. However, this time, we use two variables, $x$ and $y$.
\begin{align}
\label{eqn:defR}
&R_{A,B,C,D}(x,y,q)\\
&= \sum_{i,j,k,\ell \ge 0}
\frac{x ^{2j + k + \ell } y ^{ i+j+k}q^{2i^2 + 6 ij + 4ik + 2i\ell + 6j^2 + 8 jk + 4 j\ell + 3k^2+3k\ell+\ell^2+Ai+Bj+Ck+D\ell}}
{\left(q^2;q^2\right)_i\left(q^2;q^2\right)_j\left(q;q\right)_k\left(q;q\right)_\ell}. \nonumber
\end{align}
As usual, we will drop the arguments $x,y,q$ when they are clear.

We have the following shifts of $R_{A,B,C,D}$:
\begin{align}
R_{A,B,C,D}(xq,y,q)&=R_{A,B+2,C+1,D+1}(x,y,q) \label{eqn:shiftRx},\\
R_{A,B,C,D}(x,yq,q)&=R_{A+1,B+1,C+1,D}(x,y,q) \label{eqn:shiftRy}.
\end{align}
